\documentclass[a4paper,12pt]{report}
\usepackage[catalan]{varioref}
\usepackage{setspace}
\usepackage[margin=2.54cm]{geometry}
\usepackage{pdfpages}
\usepackage[utf8]{inputenc}
\usepackage[catalan]{babel}
\usepackage{graphicx,subcaption}
\usepackage{graphics}
\usepackage{lscape}
\usepackage{pdflscape}
\usepackage{float}
\usepackage{textcomp}
\usepackage{amsmath}
\usepackage{hyperref}
\usepackage{subcaption}
\usepackage{pgfplots}
\usepackage{tikz}
\usepackage{fancyvrb}
\usepackage{parskip}
\usepackage{changepage}
\usepackage{enumitem}
\usepackage{tcolorbox}
\usepackage[all]{hypcap}
\usepackage{xcolor}
\usepackage{listings}
\definecolor{green}{HTML}{228B22}
\definecolor{orange}{HTML}{FFC107}
\usepackage{color}
\definecolor{dkgreen}{rgb}{0,0.6,0}
\definecolor{gray}{rgb}{0.5,0.5,0.5}
\definecolor{mauve}{rgb}{0.58,0,0.82}
\lstset{escapeinside={<@}{@>}}

\hypersetup{
    colorlinks,
    citecolor=black,
    filecolor=black,
    linkcolor=black,
    urlcolor=blue
}


\lstset{frame=tb,
    language=python,
    aboveskip=3mm,
    belowskip=3mm,
    showstringspaces=false,
    columns=flexible,
    basicstyle={\small\ttfamily},
    numbers=none,
    numberstyle=\tiny\color{gray},
    keywordstyle=\color{blue},
    commentstyle=\color{dkgreen},
    stringstyle=\color{mauve},
    breaklines=true,
    breakatwhitespace=true, tabsize=3
}
\title{
	\begin{center}
	\vspace{3cm}
	\includegraphics[width=11cm, height=3cm]{images/Logo-uoc.png}
	\end{center}
	\begin{center}
	\line(1,0){400}
	\end{center}		
	TIPOLOGIA I CICLE DE VIDA DE LES DADES\\
	\vspace{2mm}
	\Large PAC1: Què son les dades i quin el seu cicle de vida?\\
	\line(1,0){400}
	\vspace{2.5cm}
	}

\author{Marc Cervera Rosell \vspace{1cm}}


\date{Semestre: febrer 2025 - juny 2025\vspace{0.5cm} \\ Màster en ciència de dades}
\onehalfspacing

\begin{document}
\thispagestyle{empty}
	\begin{titlepage}
		\maketitle
		\thispagestyle{empty}
	\end{titlepage}
	\cleardoublepage
	\newpage

\thispagestyle{empty}
\tableofcontents
\thispagestyle{empty}
\newpage
\pagenumbering{arabic}
%\thispagestyle{empty}
\section*{Exercici 1}
\addcontentsline{toc}{section}{Exercici 1}
\subsection*{Pregunta 1}
\addcontentsline{toc}{subsection}{Pregunta 1}
Per aplicar la tècnica de normalització caldria, en primer lloc, seleccionar aquelles variables numèriques a les quals volem aplicar la tècnica i després ajustar-les a una escala més reduïda, per tal de, fer-les més comprensibles. Típicament, les escales són entre -1.0 i 1.0 o entre 0.0 i 1.0.\\
Per aplicar la tècnica de discretització cal, com en el cas anterior, seleccionar aquelles variables numèriques a les quals volem aplicar la tècnica. Seguidament, cal substituir els valors numèrics per una sèrie d'etiquetes les quals poden ser conceptuals o intervals. Aquesta tècnica permet consolidar un possible criteri d'agrupació de dades. Per exemple, un \textit{dataset} amb els resultats dels balanços d'empreses. Es podria substituir els resultats dels balanços per 'positiu', 'negatiu' i 'neutre' i, posteriorment, agrupar pel resultat del balanç.\\
Aquestes tècniques són molt importants per als models de predicció, ja que permeten que les dades siguin més senzilles d'interpretar i això comporta que el rendiment del model sigui millor i que aquest tingui la capacitat de trobar patrons més fàcilment, és a dir, el model es fa més senzill.
\subsection*{Pregunta 2}
\addcontentsline{toc}{subsection}{Pregunta 2}

\subsection*{Pregunta 3}
\addcontentsline{toc}{subsection}{Pregunta 3}

\subsection*{Pregunta 4}
\addcontentsline{toc}{subsection}{Pregunta 4}

\subsection*{Pregunta 5}
\addcontentsline{toc}{subsection}{Pregunta 5}

\subsection*{Pregunta 6}
\addcontentsline{toc}{subsection}{Pregunta 6}

\subsection*{Pregunta 7}
\addcontentsline{toc}{subsection}{Pregunta 7}

\end{document}
